\chapter{Certification Exam Mapping}
\index{SnowPro certification}\index{exam mapping}\index{study strategy}%
This appendix maps the two Snowflake certifications this book serves---\textbf{SnowPro Core}
and \textbf{SnowPro Advanced: Data Engineer}---to the chapters and labs that teach each exam
objective, then lays out the study strategy and a self-assessment checklist. The official exam
guides are the authoritative blueprints and supersede anything here \cite{snowproExamGuide};
weights below reflect the published domain weightings at the time of writing. Chapter~24 is the
book's guided final review and assumes you have worked the mapping below.

\section{Which Exam, and When}
\index{SnowPro Core}\index{SnowPro Advanced}
\begin{itemize}
    \item \textbf{SnowPro Core} is the breadth exam: architecture, security, performance
    concepts, loading, transformation, and protection/sharing across the whole platform. It is
    a prerequisite for the Advanced track and is passable around \textbf{Chapter~12}, once
    Parts I--III are complete.
    \item \textbf{SnowPro Advanced: Data Engineer} is the depth exam for this book's audience:
    data movement, performance tuning, storage/protection, and security/governance, weighted
    30/30/20/20. Attempt it after completing the full 24 weeks and the Chapter~24 timed
    rehearsal.
\end{itemize}

\section{SnowPro Core: Domain-to-Chapter Map}
\index{SnowPro Core!domain map}
\begin{center}
\footnotesize
\begin{tabular}{@{}p{5.2cm}p{1.4cm}p{7.2cm}@{}}
\toprule
\textbf{Exam domain} & \textbf{Weight} & \textbf{Book coverage} \\
\midrule
Snowflake Data Cloud features \& architecture & 25\% & Ch.~1 (three layers, micro-partitions), Ch.~2 (warehouses), Ch.~0 (editions, account model), Ch.~21 (hybrid tables) \\
Account access \& security & 20\% & Ch.~17 (RBAC hierarchy, MFA, network policies, key-pair auth), Ch.~18 (masking, row access, tags) \\
Performance concepts & 15\% & Ch.~2 (sizing, scaling, caching), Ch.~3 (clustering, Search Optimization, materialized views), Ch.~10 (Query Profile) \\
Data loading \& unloading & 10\% & Ch.~5 (stages, file formats, \texttt{COPY INTO}), Ch.~6 (Snowpipe) \\
Data transformation & 20\% & Ch.~7 (streams), Ch.~8 (tasks), Ch.~9 (semi-structured), Ch.~11 (dynamic tables), Ch.~12 (window functions, dimensional design), Ch.~13--15 (Snowpark, UDFs, procedures) \\
Data protection \& data sharing & 10\% & Ch.~4 (Time Travel, Fail-safe, cloning), Ch.~19 (shares, reader accounts, Marketplace), Ch.~20 (clean rooms) \\
\bottomrule
\end{tabular}
\end{center}

\noindent Chapters 16 and 22--24 (external functions, Cortex, DataOps, exam prep) extend past
Core scope but reinforce it; Chapters 13--16 supply the programmability questions' context even
where Core stays SQL-level.

\section{SnowPro Advanced: Data Engineer---Domain-to-Chapter Map}
\index{SnowPro Advanced!domain map}
\begin{center}
\footnotesize
\begin{tabular}{@{}p{4.6cm}p{1.4cm}p{7.8cm}@{}}
\toprule
\textbf{Exam domain} & \textbf{Weight} & \textbf{Book coverage} \\
\midrule
Data movement & 30\% & Ch.~5 (stages, storage integrations, \texttt{COPY INTO} semantics), Ch.~6 (Snowpipe, streaming, offset tokens), Ch.~7 (streams, CDC), Ch.~8 (tasks, DAGs), Ch.~11 (dynamic tables), Ch.~9 (semi-structured ingestion) \\
Performance tuning & 30\% & Ch.~1 (pruning, clustering depth), Ch.~2 (sizing, multi-cluster policies), Ch.~3 (Search Optimization, materialized views), Ch.~10 (profile reading, spill, joins), Ch.~12 (SQL shape) \\
Storage \& data protection & 20\% & Ch.~4 (Time Travel, Fail-safe, cloning, \texttt{UNDROP}), Ch.~3 (table types and retention), Ch.~21 (hybrid tables) \\
Security \& governance & 20\% & Ch.~17 (RBAC, networking, secrets), Ch.~18 (masking, row access, tags), Ch.~19 (sharing, reader accounts), Ch.~20 (clean rooms, erasure) \\
\bottomrule
\end{tabular}
\end{center}

\noindent The Snowpark chapters (13--15) and DataOps chapter (23) appear throughout Advanced
scenario stems---expect questions that embed Snowpark client behavior, procedure rights models,
and pipeline-deployment judgment inside the four weighted domains.

\section{Study Strategy}
\index{study strategy}
Chapter~24 details the rehearsal protocol; the strategy in brief:
\begin{enumerate}
    \item \textbf{Allocate hours by weight.} For Advanced, 30/30/20/20 means movement and
    performance earn 60\% of study time. For Core, architecture and transformation dominate.
    \item \textbf{Rehearse failure modes, not happy paths.} Scenario questions live where
    things break: a stale stream, a wrong event prefix, a third task retry, an expired
    retention window. For every feature, be able to state what fails and what the platform does
    next.
    \item \textbf{Memorize the numbers.} Retention windows (Time Travel by edition, Fail-safe's
    seven days, stream staleness), load-metadata duration (64 days), scaling policy behavior,
    and billing units are high-frequency fact anchors.
    \item \textbf{Practice elimination.} Read the final question sentence first, extract
    constraints (edition, latency, budget, semantics), and eliminate options that violate one
    before judging survivors (Ch.~24).
    \item \textbf{Measure, then remediate hands-on.} Score timed mocks \emph{by domain}; every
    miss becomes an error-log entry (topic, misread signal, correct rule, book section), and
    every remediable topic gets its lab re-run, not just re-read (Ch.~24).
    \item \textbf{Book the exam on evidence.} Two consecutive timed attempts at 80\%+ with no
    domain under 70\% is the readiness bar used in this book.
\end{enumerate}

\section{Self-Assessment Checklist}
\index{self-assessment}
Answer each item from memory; any hesitation marks a chapter to revisit. Appendix~F provides
the corresponding interview-style drills.

\textbf{Architecture and compute (Ch.~0--3).}
\begin{itemize}
    \item[$\square$] I can name the three architecture layers and what each bills for.
    \item[$\square$] I can explain micro-partitions, pruning, and why a clustering key is not an index.
    \item[$\square$] I can choose a warehouse size and scaling policy for a described workload and defend it in credits.
    \item[$\square$] I can distinguish result cache, local disk cache, and metadata cache.
    \item[$\square$] I can state when transient/temporary tables and Search Optimization are correct.
\end{itemize}

\textbf{Data movement (Ch.~5--9, 11).}
\begin{itemize}
    \item[$\square$] I can build a stage + storage integration with no embedded credentials.
    \item[$\square$] I know \texttt{COPY INTO} idempotency rules, \texttt{ON\_ERROR} modes, and \texttt{VALIDATION\_MODE}.
    \item[$\square$] I can contrast Snowpipe, Snowpipe Streaming, and bulk COPY on latency, billing, and delivery guarantees.
    \item[$\square$] I can explain stream offsets, \texttt{METADATA\$ISUPDATE}, and the staleness limit.
    \item[$\square$] I can design a task DAG with gating, timeouts, and failure containment.
    \item[$\square$] I can set a dynamic-table lag budget and detect incremental $\rightarrow$ full refresh fallbacks.
\end{itemize}

\textbf{Transformation and programmability (Ch.~10--16).}
\begin{itemize}
    \item[$\square$] I can read a Query Profile: dominant operator, partitions scanned/total, spill, join pathology.
    \item[$\square$] I can use \texttt{QUALIFY}, GROUPING SETS, and PIVOT, and design a star schema with declared grain.
    \item[$\square$] I can explain Snowpark lazy evaluation and verify pushdown with \texttt{explain()}.
    \item[$\square$] I can choose among SQL UDF, vectorized Python UDF, and UDTF, and state the Anaconda channel's role.
    \item[$\square$] I can contrast owner's-rights and caller's-rights procedures with a use case for each.
    \item[$\square$] I can wire an external function (API integration + remote service) and explain its cost shape.
\end{itemize}

\textbf{Governance, protection, and sharing (Ch.~4, 17--21).}
\begin{itemize}
    \item[$\square$] I can contrast Time Travel, Fail-safe, and zero-copy cloning, including retention by object type.
    \item[$\square$] I can design the access-role/functional-role hierarchy for a three-team account.
    \item[$\square$] I can implement tag-based masking and a row access policy, and test both as a non-owner role.
    \item[$\square$] I can explain crypto-shredding for erasure across clones and Fail-safe.
    \item[$\square$] I can choose between a direct share, a reader account, and a Marketplace listing, and say who pays compute.
    \item[$\square$] I can state what hybrid tables enforce and when a standard table is still correct.
\end{itemize}

\textbf{Operations and AI (Ch.~22--23).}
\begin{itemize}
    \item[$\square$] I can set account and warehouse resource monitors with notify/suspend triggers.
    \item[$\square$] I can attribute credits by team using tags and \texttt{ACCOUNT\_USAGE}.
    \item[$\square$] I can describe the Terraform + dbt + schemachange deployment flow and its CI authentication.
    \item[$\square$] I know what DMFs monitor and where their results land.
    \item[$\square$] I can name the Cortex function families, the \texttt{VECTOR} type, and how Cortex spend is capped.
\end{itemize}

\begin{notebox}
Checkboxes are claims, not proof. For every box you tick, be ready to produce the lab evidence:
the Query Profile screenshot, the masking-policy negative test, the resource-monitor trigger
history. That evidence doubles as the Chapter~24 portfolio.
\end{notebox}
